\begin{helper}[FiberAndDegreeMixedLiftedIntersectionConcentration]
\label{helper:FiberAndDegreeMixedLiftedIntersectionConcentration}
Let $p$ be the edge probability and let $m$ be the reduced vertex count.
With high probability, the two-bite sample has the property that for every red base vertex $r$ and every blue base vertex $b$, the intersection of their lifted neighborhoods has size at most $100 \log^3 n$.
\end{helper}

\begin{proof}
Write \(L=\log n\), \(M=m(n)=\noderef{TwoBiteNaturalReducedVertexCount}(n)\), and
\[
p=\noderef{TwoBiteEdgeProbability}(1/2,n)=\frac12\sqrt{\frac Ln}.
\]
By \(\noderef{TwoBiteNaturalReducedVertexCount}\), together with
\(\noderef{TwoBiteReducedVertexCount}\) and \(\noderef{TwoBiteStretch}\), for
all sufficiently large \(n\),
\[
M=\left\lfloor\frac n{L^2}\right\rfloor,\qquad
\frac12\frac n{L^2}\le M\le\frac n{L^2},\qquad n\le M^2.
\tag{1}
\]
The eventual support and total-mass facts for this rounded size are supplied
by \(\noderef{TwoBiteNaturalTotalMassOneEventually}\), and the probability
bounds \(0\le p\le1\) are supplied eventually by
\(\noderef{OppositeProjectionEdgeProbBounds}\).  Discarding finitely many
values of \(n\) does not affect the \(\noderef{WithHighProbability}\)
conclusion.

Let \(\mathcal D_n\) be the event that every red base degree and every blue
base degree is at most \(2pM\).  This event holds with high probability by
\(\noderef{FiberAndDegreeBaseDegreeConcentration}\), applied with tolerance
\(1\): the conclusion of that node says that all red and blue base degrees
are within relative error \(1\) of \(pM\), and hence each is at most
\((1+1)pM=2pM\).  It remains to prove that, conditional on
\(\mathcal D_n\), the probability that some mixed lifted intersection
exceeds \(100L^3\) tends to \(0\).

We now spell out the conditional distribution of the injection.  By
\(\noderef{TwoBiteSample}\), an outcome is a triple
\[
\omega=(G_R,G_B,\phi),
\]
where \(G_R=\noderef{TwoBiteRedGraph}(\omega)\) and
\(G_B=\noderef{TwoBiteBlueGraph}(\omega)\) are the two base graphs on
\(\operatorname{Fin}(M)\), and \(\phi\) is an injection of the \(n\) final
vertices into the \(M^2\) cells
\(\operatorname{Fin}(M)\times\operatorname{Fin}(M)\).  By
\(\noderef{TwoBiteEventProbability}\), probabilities are finite sums of
\(\noderef{TwoBiteSampleWeight}\) over the relevant triples.  The summands
are nonnegative by \(\noderef{TwoBiteSampleWeightNonnegative}\).  The weight
factorization in \(\noderef{TwoBiteSampleWeight}\) has the form
\[
\noderef{GnpGraphWeight}(p,G_R)\,
\noderef{GnpGraphWeight}(p,G_B)\,
\noderef{UniformInjectionWeight}(n,M).
\tag{2}
\]
When one marginalizes an unused base graph, the corresponding graph-weight
sum is \(1\) by \(\noderef{GnpGraphWeightSumOne}\).  When both base graphs
are fixed, the two graph-weight factors in (2) are constants in the
conditional numerator and denominator, so they cancel.  By (1) and
\(\noderef{TwoBiteNaturalTotalMassOneEventually}\), the injection support is
nonempty; hence \(\noderef{UniformInjectionWeight}(n,M)\) is the reciprocal
of the number of injections and its total mass over injections is \(1\).
Therefore, conditional on the two fixed base graphs, \(\phi\) is a uniform
ordered sample of \(n\) distinct cells from the \(M^2\) product cells.

Fix base graphs satisfying \(\mathcal D_n\), and fix a red vertex
\(r\in\operatorname{Fin}(M)\) and a blue vertex \(b\in\operatorname{Fin}(M)\).
Let
\[
d_R(r)=|N_{G_R}(r)|,\qquad d_B(b)=|N_{G_B}(b)|.
\]
On \(\mathcal D_n\), both \(d_R(r)\) and \(d_B(b)\) are at most \(2pM\).
By the definitions of \(\noderef{TwoBiteLiftedNeighborhood}\),
\(\noderef{TwoBiteRedProjection}\), \(\noderef{TwoBiteBlueProjection}\), and
\(\noderef{TwoBiteEmbedding}\), a final vertex lies in the mixed intersection
\[
\noderef{TwoBiteLiftedNeighborhood}(\omega,\operatorname{inl}(r))
\cap
\noderef{TwoBiteLiftedNeighborhood}(\omega,\operatorname{inr}(b))
\]
exactly when its injected product cell lies in the marked rectangle
\[
K_{r,b}=N_{G_R}(r)\times N_{G_B}(b).
\]
Thus
\[
|K_{r,b}|=d_R(r)d_B(b)\le (2pM)(2pM)=4p^2M^2.
\tag{3}
\]

Let \(Y_{r,b}\) be the size of this mixed lifted intersection.  Conditional
on \(G_R,G_B\), the previous paragraph identifies \(Y_{r,b}\) as a
hypergeometric, without-replacement count from population \(M^2\), marked
set \(K_{r,b}\), and sample size \(n\).  Its conditional mean is
\[
\mu_{r,b}=n\frac{|K_{r,b}|}{M^2}\le 4p^2n=L.
\tag{4}
\]
This is precisely the without-replacement model controlled by
\(\noderef{HypergeometricMgfComparison}\); for the present large upper tail,
we use the following direct counting consequence.  Put
\[
T=\lceil 100L^3\rceil.
\]
If \(Y_{r,b}\ge T\), then some \(T\)-element subset of \(K_{r,b}\) is
contained in the image of the uniform injection.  For a fixed \(T\)-subset
of \(K_{r,b}\), the probability that all its cells appear in the image is
\[
\frac{(n)_T}{(M^2)_T}\le\left(\frac n{M^2}\right)^T.
\]
Taking a union bound over the \(\binom{|K_{r,b}|}{T}\) possible subsets and
using \(\binom{K}{T}\le(eK/T)^T\), we obtain
\[
\Pr(Y_{r,b}\ge T\mid G_R,G_B)
\le
\binom{|K_{r,b}|}{T}\left(\frac n{M^2}\right)^T
\le
\left(\frac{e\mu_{r,b}}{T}\right)^T.
\tag{5}
\]
By (4) and \(T\ge100L^3\),
\[
\left(\frac{e\mu_{r,b}}{T}\right)^T
\le
\left(\frac{e}{100L^2}\right)^T
\le \exp(-cL^3\log L)
\tag{6}
\]
for some absolute constant \(c>0\) and all sufficiently large \(n\).  Since
the event \(Y_{r,b}>100L^3\) implies \(Y_{r,b}\ge T\), (5)--(6) bound the
failure probability for this fixed mixed pair.

There are \(M^2\) choices of the mixed pair \((r,b)\), and by (1) we have
\(M^2\le n^2=\exp(2L)\).  On \(\mathcal D_n\), the union bound therefore
gives
\[
\Pr\!\left(\exists r,b:\,
Y_{r,b}>100L^3 \,\middle|\, G_R,G_B\right)
\le
\exp(2L)\exp(-cL^3\log L)=o(1).
\tag{7}
\]
Finally \(\Pr(\mathcal D_n)\to1\), so the unconditional failure probability
is at most the \(o(1)\) failure probability of \(\mathcal D_n\) plus the
\(o(1)\) bound in (7).  Hence, with high probability, every mixed red-blue
lifted-neighborhood intersection has size at most \(100(\log n)^3\), as
claimed.
\end{proof}
